% ===== Block F: PDF p.27–30 | 原文 §3 Cohomological field theories 之 §3.1 Axioms（两标题之后） | 唯一编号公式 3.1–3.24 =====
% [NOTE: 本块从 p.27 的 “3. Cohomological field theories” 标题之后（即 p.27 页顶第一段）开始；该级标题由 Lead 在 main.tex 中撰写。原文的 “3.1. Axioms” 标题排在 p.27 中部（(3.3) 与 Kontsevich--Manin 起源段之后），按 Lead 裁决由本块就地写入 \subsection，故本块不写 \section。]
% [NOTE: 原文 p.28 的 (3.9)-(3.12) 与 (3.15) 在原书中是图示（曲面粘合图与配边图）而非代数式；按本项目“不画图”的统一规定，本块以文字描述与拓扑标签如实转述，编号仍与原文一致。]
% [原文 p.27]

Witten 的关联函数所满足的 Virasoro 约束，给出了计算全部 $\psi$ 类相交数的一种递归方法。这些约束所依托的主要几何性质，是 $\Mc_{g,n}$ 的递归结构。把 Witten 的关联函数看作单位元与 $\psi$ 类的相交，我们就能把这一递归结构纯粹用上同调的语言重新表述。单位元 $\mathbf{1}_{g,n}\in H^{0}(\Mc_{g,n},\Q)$ 在一切重言映射的拉回下都是稳定的，即

\begin{align}
  \rho^{*}\mathbf{1}_{g,n} &= \mathbf{1}_{g-1,n+2}\,, & \rho &:\Mc_{g-1,n+2}\longrightarrow\Mc_{g,n}\,, \tag{3.1}\\
  \sigma^{*}\mathbf{1}_{g,n} &= \mathbf{1}_{g_{1},1+|I_{1}|}\otimes\mathbf{1}_{g_{2},1+|I_{2}|}\,, & \sigma &:\Mc_{g_{1},1+|I_{1}|}\times\Mc_{g_{2},1+|I_{2}|}\longrightarrow\Mc_{g,n}\,, \tag{3.2}\\
  \pi^{*}\mathbf{1}_{g,n} &= \mathbf{1}_{g,n+1}\,, & \pi &:\Mc_{g,n+1}\longrightarrow\Mc_{g,n}\,. \tag{3.3}
\end{align}

% [Tutor Note: 原文把单位类排成粗体的 $\mathbf{1}_{g,n}$，以区别于数字 $1$、基向量 $e_1$ 以及 (3.4) 中出现的下标 $1$；本译注本沿用这一粗体约定。核对方法：$H^{0}(\Mc_{g,n},\Q)$ 是一维的，故 $\rho^{*}\mathbf{1}_{g,n}$ 必为 $\mathbf{1}_{g,n}$ 的常数倍，而 (3.1) 断言该常数为 $1$，也就是说单位类就是基本类 $[\Mc_{g,n}]$，它在所有重言映射下都不被“压掉”。]
% [Tutor Note: (3.2) 右端两个因子 $\mathbf{1}_{g_{1},1+|I_{1}|}$、$\mathbf{1}_{g_{2},1+|I_{2}|}$ 各比 $\mathbf{1}_{g_{1},|I_{1}|}$、$\mathbf{1}_{g_{2},|I_{2}|}$ 多出一个下标 $1$：粘合时每条粘合边在两侧各贡献一个标记点。核对方法：取 $n=0$，则 $I_{1}=I_{2}=\varnothing$，(3.2) 化为 $\sigma^{*}\mathbf{1}_{g,0}=\mathbf{1}_{g_{1},1}\otimes\mathbf{1}_{g_{2},1}$，正是“两块曲线各带一个粘合点”的边界除子定义。]

前两个等式可以解释为 2D 拓扑场论（topological field theory，简称 TQFT）中局部性公理的上同调版本。受 TQFT 的启发，我们下面基于 $\Mc_{g,n}$ 的上同调来定义它们的上同调类比。原始定义由 M. Kontsevich 与 Y. Manin 在 90 年代中期给出 [KM94]，那是将拓扑弦论公理化的首次尝试，并与 B. Dubrovin 关于 2D TQFT 几何的开创性工作 [Dub96] 有深刻的联系。

\subsection{公理（Axioms）}

固定一个有限维 $\Q$-向量空间 $V$，称为\emph{相空间}（phase space），并在其上配备一个非退化配对 $\eta:V\times V\to\Q$。为方便起见，我们在 $V$ 的一个固定基 $(e_{1},\dots,e_{r})$ 下工作。记 $(\eta_{\mu,\nu})$ 为该配对的矩阵元素，而 $(\eta^{\mu,\nu})$ 为其逆矩阵。

\begin{definition}[3.1]
$(V,\eta)$ 上的一个\emph{上同调场论}（cohomological field theory）由一族线性映射 $\Omega=(\Omega_{g,n})_{2g-2+n>0}$ 组成：
\begin{equation}
  \Omega_{g,n}:V^{\otimes n}\longrightarrow H^{2\bullet}(\Mc_{g,n},\Q)\,,\qquad
  \Omega_{g,n}(e_{\mu_{1}}\otimes\cdots\otimes e_{\mu_{n}})=\Omega_{g;\mu_{1},\dots,\mu_{n}}\,,
  \tag{3.4}
\end{equation}
且满足下列公理。
% [Tutor Note: 指标条件 $2g-2+n>0$ 是稳定曲线的存在条件，它恰好排除 $(g,n)=(0,0),(0,1),(0,2),(1,0)$ 四种使 $\Mc_{g,n}$ 不稳定（或不存在）的情形。核对方法：$(0,3)$ 满足 $2\cdot0-2+3=1>0$，故 $(0,3)$ 是允许的，这正是 (3.8) 中 $\Omega_{0;\mu,\nu,1}$ 有意义的前提。]
\begin{enumerate}[label=\roman*)]
  \item \emph{对称性}。每个 $\Omega_{g,n}$ 都是 $S_{n}$-不变的，其中对称群 $S_{n}$ 的作用同时置换 $\Mc_{g,n}$ 的标记点与 $V^{\otimes n}$ 的各分量。
  \item \emph{粘合}。考虑粘合映射
  \begin{align}
    \rho &:\Mc_{g-1,n+2}\longrightarrow\Mc_{g,n}\,, \nonumber\\
    \sigma &:\Mc_{g_{1},1+|I_{1}|}\times\Mc_{g_{2},1+|I_{2}|}\longrightarrow\Mc_{g,n}\,,\qquad g_{1}+g_{2}=g,\ I_{1}\sqcup I_{2}=\{1,\dots,n\}\,,
    \tag{3.5}
  \end{align}
  我们要求
  \begin{align}
    \rho^{*}\Omega_{g;\mu_{1},\dots,\mu_{n}} &= \eta^{\alpha,\beta}\,\Omega_{g-1;\alpha,\beta,\mu_{1},\dots,\mu_{n}}\,, \nonumber\\
    \sigma^{*}\Omega_{g;\mu_{1},\dots,\mu_{n}} &= \eta^{\alpha,\beta}\,\Omega_{g_{1};\alpha,\mu_{I_{1}}}\otimes\Omega_{g_{2};\beta,\mu_{I_{2}}}\,.
    \tag{3.6}
  \end{align}
% [Tutor Note: (3.6) 用逆配对 $\eta^{\alpha,\beta}$（上标）缩并，与 (3.5) 中 $\rho$、$\sigma$ 的定义域一一对应：$\rho$ 一侧多出两个指标 $\alpha,\beta$，$\sigma$ 一侧两块各多出一个指标 $\alpha$、$\beta$，再用 $\eta^{\alpha,\beta}$ 把这一对指标缩并掉。核对方法：取平凡 CohFT $\Omega_{g,n}=\mathbf{1}_{g,n}$，此时 $\eta^{\alpha,\beta}\mathbf{1}\otimes\mathbf{1}=1$，(3.6) 退化为 (3.2)，可见 (3.1)-(3.3) 是 (3.6) 的特例。]
  \item[] 若该向量空间还带有一个特殊的非零元素（不失一般性，可以取为 $e_{1}$），我们还可以要求第三条公理。
% [原文 p.28]
  \item \emph{单位}。考虑遗忘映射
  \begin{equation}
    \pi:\Mc_{g,n+1}\longrightarrow\Mc_{g,n}\,,
    \tag{3.7}
  \end{equation}
  我们要求
  \begin{equation}
    \pi^{*}\Omega_{g;\mu_{1},\dots,\mu_{n}}=\Omega_{g;\mu_{1},\dots,\mu_{n},1}
    \qquad\mbox{以及}\qquad
    \Omega_{0;\mu,\nu,1}=\eta_{\mu,\nu}\,.
    \tag{3.8}
  \end{equation}
\end{enumerate}
\end{definition}

此时称 $\Omega$ 为\emph{带单位的上同调场论}；该特殊元素称为\emph{单位}或\emph{真空}。

这些公理可以用图形表示如下。

\begin{equation}
  \parbox{0.9\textwidth}{\centering\small\kaishu 〔原文此处为图示：左端是亏格为 $g$、leg标 $\mu_{1},\dots,\mu_{n}$ 的曲面上的 $\Omega_{g}$；经 $\rho^{*}$ 拉回后，右端为亏格 $g-1$ 曲面上的 $\Omega_{g-1}$，并多出一对以 $\alpha$（上方）、$\beta$（下方）标记的leg，这对leg以 $\eta$ 缩并。〕}
  \tag{3.9}
\end{equation}

\begin{equation}
  \parbox{0.9\textwidth}{\centering\small\kaishu 〔原文此处为图示：左端同 (3.9) 中的 $\Omega_{g}$；经 $\sigma^{*}$ 拉回后，右端分裂为 $\Omega_{g_{1}}$ 与 $\Omega_{g_{2}}$ 两块，二者以一条标 $\eta$ 的边相连（在 $\Omega_{g_{1}}$ 一侧标 $\alpha$，在 $\Omega_{g_{2}}$ 一侧标 $\beta$），两块余下的leg分别标 $\mu_{I_{1}}$ 与 $\mu_{I_{2}}$。〕}
  \tag{3.10}
\end{equation}

\begin{equation}
  \parbox{0.9\textwidth}{\centering\small\kaishu 〔原文此处为图示：经 $\pi^{*}$，$\Omega_{g}$ 变为其自身，只是多出一条标 $1$ 的leg。〕}
  \tag{3.11}
\end{equation}

\begin{equation}
  \parbox{0.9\textwidth}{\centering\small\kaishu 〔原文此处为图示：带leg $\mu,\nu$ 且另有一条标 $1$ 的leg的亏格 $0$ 曲面上的 $\Omega_{0}$，等于只带leg $\mu,\nu$ 的“帽”；右端的帽即 $\eta_{\mu,\nu}$。〕}
  \tag{3.12}
\end{equation}

% [Tutor Note: (3.9)-(3.12) 在原书中是四幅图示而非代数式，本译注本按“不画图”的统一规定改用文字描述，编号照原文保留。其代数内容依次是 (3.6) 第一式、(3.6) 第二式、(3.8) 第一式与 (3.8) 第二式：换言之这四条“公式”只是把公理 (ii)、(iii) 画成了曲面粘合图，本身不引入新的等式。核对方法：在 (3.9) 中数亏格与标记数，$g\to g-1$、$n\to n+2$，与 $\rho$ 的定义域 $\Mc_{g-1,n+2}$ 吻合。]

一个上同调场论（简称 CohFT）决定了 $V$ 上的一个乘积 $\star$，称为\emph{量子积}：
\begin{equation}
  e_{\mu}\star e_{\nu}=\Omega_{0;\mu,\nu,\alpha}\,\eta^{\alpha,\beta}\,e_{\beta}\,.
  \tag{3.13}
\end{equation}

$\star$ 的交换性与结合性分别由 (i) 与 (ii) 推出。若该 CohFT 带单位，则量子积是单位的，由 (iii) 知 $e_{1}\in V$ 就是单位元。

% [Tutor Note: (3.13) 之所以是 $V$ 上的乘积（而不取值于某个形式幂级数环），是因为 $\Mc_{0,3}$ 是单点，$H^{2\bullet}(\text{pt},\Q)=\Q$，故 $\Omega_{0;\mu,\nu,\alpha}$ 是纯数、经 $\eta^{\alpha,\beta}$ 升降指标后仍在 $V$ 中。核对方法：由 (3.13) 与配对可得 $\eta(e_{\mu}\star e_{\nu},e_{\rho})=\Omega_{0;\mu,\nu,\rho}$；交换性即 $\Omega_{0;\mu,\nu,\rho}$ 对前两个指标对称，而这正是公理 (i) 在 $n=3$ 时的内容。]

\begin{exercise}[3.1]
证明 $(V,\eta,\star)$ 构成一个 Frobenius 代数，即它满足
\begin{equation}
  \eta(v_{1}\star v_{2},v_{3})=\eta(v_{1},v_{2}\star v_{3})\,.
  \tag{3.14}
\end{equation}
\end{exercise}

% [Tutor Note: 提示（不写进译文）：由 (3.13) 有 $\eta(v_{1}\star v_{2},v_{3})=\Omega_{0;\mu\nu\rho}v_{1}^{\mu}v_{2}^{\nu}v_{3}^{\rho}$，故 (3.14) 等价于 $\Omega_{0;\mu\nu\rho}$ 对三个指标全对称；对称性来自公理 (i)（$S_{3}$ 同时置换三个标记点），而结合性来自公理 (ii) 在 $\Mc_{0,4}$ 上的两种粘合方式给出的同一个等式。核对方法：$g=0$、$n=3$ 时没有可粘合的边界，须升到 $\Mc_{0,4}$ 才能看出结合性。]

一个（带单位 $e$ 的）Frobenius 代数通过下列对应关系等价于一个 2D 拓扑场论 $Z$：$Z(S^{1})=V$ 为圆周上的态 Hilbert 空间，而
\begin{align}
  Z\big(\mbox{〔两进零出〕}\big) &= \eta:V\otimes V\longrightarrow\Q\,, \nonumber\\
  Z\big(\mbox{〔两进一出〕}\big) &= \star:V\otimes V\longrightarrow V\,, \nonumber\\
  Z\big(\mbox{〔零进一出〕}\big) &= e:\Q\longrightarrow V
  \tag{3.15}
\end{align}
为态射。任意亏格为 $g$、连接 $n$ 个初态与 $m$ 个末态的曲面 $\Sigma_{g,n,m}$ 的配分函数 $Z(\Sigma_{g,n,m})$，都可以用上述取值借助 TFT 的性质重构出来。

% [Tutor Note: (3.15) 三行在原书中同样是图示：依次是“两条leg进、无leg出”的帽（给出 $\eta$）、“两进一出”的裤子（给出 $\star$）、“无进一出”的圆盘（给出 $e$）；本译注本用进、出leg数作为拓扑标签取代图形。核对方法：把裤子（两进一出）的三条leg按两种方式配对切开，得到的正是 (3.14) 的配对不变性，这就是“Frobenius 代数等价于 2D TFT”的几何来源。]

% [原文 p.29]

与任何 CohFT $\Omega$ 相伴，我们还得到一族有理数，称为\emph{CohFT 关联函数}（或\emph{祖先不变量}，ancestor invariants），它们定义为
\begin{equation}
  \left\langle\tau_{\mu_{1},d_{1}}\cdots\tau_{\mu_{n},d_{n}}\right\rangle^{\Omega}_{g}
  =\int_{\Mc_{g,n}}\Omega_{g;\mu_{1},\dots,\mu_{n}}\prod_{i=1}^{n}\psi_{i}^{d_{i}}\,.
  \tag{3.16}
\end{equation}

注意，由于次数上的理由，$\sum_{i=1}^{n}d_{i}\le 3g-3+n$。

% [Tutor Note: (3.16) 中 $\psi_{i}$ 是第 $i$ 个标记点处的第一 Chern 类（见 §2.3），其实次数为 $2$；而 $\dim\Mc_{g,n}=3g-3+n$，故 $\sum_{i}d_{i}>3g-3+n$ 时积分为零，这就是原文“次数上的理由”。核对方法：$g=0,n=3$ 时 $\dim=0$，只允许 $d_{1}=d_{2}=d_{3}=0$，得到 $\langle\tau_{\mu_{1},0}\tau_{\mu_{2},0}\tau_{\mu_{3},0}\rangle=\Omega_{0;\mu_{1}\mu_{2}\mu_{3}}$，与 (3.13) 的定义相容。]

\begin{example}[3.2]
这里给出一维情形的一些 CohFT 例子。取 $V=\Q\cdot e_{1}$ 且 $\eta(e_{1},e_{1})=1$。此时我们对 $\Omega_{g,n}(e_{1}^{\otimes n})=\Omega_{g;1,\dots,1}$ 采用更简单的记号 $\Omega_{g,n}$。
% [Tutor Note: 一维情形下 $\eta(e_{1},e_{1})=1$，故指标的升降是恒等映射，$\Omega_{g;1,\dots,1}$ 与 $\Omega_{g,n}$ 可以不分；这是后文反复使用的约定。核对方法：此时 (3.16) 就是 $\Mc_{g,n}$ 上 $\psi$ 类的标准相交数，而 (3.13) 给出的量子积满足 $e_{1}\star e_{1}=\Omega_{0;111}e_{1}$。]
\begin{itemize}
  \item 令 $\Omega_{g,n}=\mathbf{1}_{g,n}$（上同调中的单位元），我们得到一个带单位 $e_{1}$、且集中在次数 $0$ 处的 CohFT。它称为\emph{平凡 CohFT}，本节开头已讨论过。相应的关联函数满足 Virasoro 约束 (2.69)，这等价于在 Airy 谱曲线 $\frac{1}{2}y^{2}-x=0$ 上的拓扑递归。
  \item 类 $\Omega_{g,n}=\exp(2\pi^{2}\kappa_{1})$ 定义一个 CohFT，由于它与双曲几何及 JT 引力的联系，有时称为\emph{Weil--Petersson CohFT}（参见 [Joh; Tur26]）。它不是带单位的 CohFT。相应的关联函数满足 Virasoro 约束（(2.69) 的 dilaton 平移版本），这等价于在正弦谱曲线（sine spectral curve）上的拓扑递归（见习题 3.2）。
% [Tutor Note: 为什么 $\exp(2\pi^{2}\kappa_{1})$ 不是“带单位”的？遗忘映射下有标准关系 $\pi^{*}\kappa_{1}=\kappa_{1}-\psi_{n+1}$，故 (3.8) 左端为 $\exp\!\big(2\pi^{2}(\kappa_{1}-\psi_{n+1})\big)$，而右端（“多一条标 $1$ 的leg”）没有这个 $\psi_{n+1}$ 因子，两者不等。另注意 $\exp$ 是类的幂级数，由 $\dim\Mc_{g,n}<\infty$ 自动截断，故 $\Omega_{g,n}$ 确实是有限和。核对方法：取 $g=0,n=4$ 计算两边在 $\Mc_{0,5}$ 上的差别。]
  \item \emph{Hodge 类} $\Omega_{g,n}=\Lambda(u)=\sum_{k=0}^{g}\lambda_{k}u^{k}$ 定义一族以 $e_{1}$ 为单位的单参数 CohFT。它作为 $\PP^{1}$ 的拓扑弦振幅的局部化公式中的顶点贡献出现。一个推广由 Hodge 类的乘积给出：
  \begin{equation}
    \Omega_{g,n}=\prod_{m=1}^{D}\Lambda(u_{m})\,,
    \tag{3.17}
  \end{equation}
  它作为 $D$ 维时空的拓扑弦振幅的局部化公式中的顶点项出现。一个特别好的情形是 $D=3$ 以及参数 $(u_{1},u_{2},u_{3})$ 满足约束
  \begin{equation}
    \frac{1}{u_{1}}+\frac{1}{u_{2}}+\frac{1}{u_{3}}=0\,.
    \tag{3.18}
  \end{equation}
  在局部化公式的语境中，该约束对应于局部 Calabi--Yau 条件 [MV02; LLZ03; OP04]（参见 [Liu]）。与 Virasoro 约束/拓扑递归的联系，目前只知道 $D=1$ 以及带有 Calabi--Yau 条件的 $D=3$ 两种情形（见习题 3.5）。
% [Tutor Note: 这里 $\lambda_{k}=c_{k}(\mathbb{E})$ 是 Hodge 丛 $\mathbb{E}$ 的 Chern 类，$k$ 只到 $g$，因为 $\mathbb{E}$ 的秩为 $g$。核对方法：$g=0$ 时 $\Lambda(u)=1$，与“$\Mc_{0,n}$ 上 Hodge 丛平凡”一致。另注意 (3.17) 是 $D$ 个 Hodge 类因子之积（不是张量积）：每个 $\Lambda(u_{m})$ 都是同一个 $\Mc_{g,n}$ 上的类，故可相乘；$D=1$ 时它退回前面的 $\Lambda(u_{1})$。]
% [Tutor Note: (3.18) 的 $\frac{1}{u_{1}}+\frac{1}{u_{2}}+\frac{1}{u_{3}}=0$ 是局部 Calabi--Yau（“中心荷抵消”）条件，其形式与三维局部曲面几何的局部化公式中三块定点贡献的权重之和为零相同。核对方法：把 $u_{m}$ 视为 $\PP^{1}$ 上定点处的等变权，该式正是三处权重的倒数之和为零，与 $c_{1}=0$ 的 Calabi--Yau 条件一致。]
  \item 在 [Nor23] 中，Norbury 定义了一个 CohFT，记作 $\Theta_{g,n}\in H^{2(2g-2+n)}(\Mc_{g,n},\Q)$，称为\emph{$\Theta$ 类}，它满足单位公理的另一个版本，即
  \begin{equation}
    \psi_{n+1}\cdot\pi^{*}\Theta_{g,n}=\Theta_{g,n+1}\,.
    \tag{3.19}
  \end{equation}
  它出现在超 JT 引力中，与 Weil--Petersson 体积的费米子部分有关 [SW20]。Norbury 猜测相应的配分函数与 KdV 层级的所谓 Br\'ezin--Gross--Witten tau 函数 [BG80; GW80] 重合，这一猜测现已在 [CGG25] 中证明。等价地，它满足 Virasoro 约束，这等价于在 Bessel 谱曲线 $\frac{1}{2}y^{2}x-1=0$ 上的拓扑递归。
% [Tutor Note: (3.19) 与 (3.8) 的差别只在多出的 $\psi_{n+1}$ 因子：普通单位公理是“多一条标 $1$ 的leg不改变类”，而 $\Theta$ 类满足的是经过 $\psi$ 修正的版本。核对次数：$\Theta_{g,n}$ 的次数为 $2(2g-2+n)$（实次数），$\psi$ 的实次数为 $2$，故 (3.19) 左端次数为 $2(2g-2+n)+2=2(2g-2+n+1)$，与右端 $2(2g-2+(n+1))$ 相等，两边自洽。]

\end{itemize}

下面列举文献中出现的一些高维 CohFT。

% [原文 p.30]

\begin{itemize}
  \item 在 [Wit92] 中，Witten 通过考虑能级为 $k$ 的 Wess--Zumino--Witten 模型（方便地重新参数化为 $k=r-2$）研究了他关于 2D 量子引力的原始工作的一个推广。这样的理论定义一个维数为 $r-1$ 的 CohFT，称为\emph{Witten $r$-旋类}（Witten $r$-spin class），其基本分量描述如下。取 $V=\bigoplus_{\mu=1}^{r-1}\Q\cdot e_{\mu}$，配对为 $\eta(e_{\mu},e_{\nu})=\delta_{\mu+\nu,r}$，单位为 $e_{1}$。Witten $r$-旋类是一个 CohFT
  \begin{equation}
    W^{r}_{g;\mu_{1},\dots,\mu_{n}}\in H^{2D^{r}_{g;\mu}}(\Mc_{g,n},\Q)
    \tag{3.20}
  \end{equation}
  其纯复次数为
  \begin{equation}
    D^{r}_{g;\mu}=\frac{(r-2)(g-1)-n+\sum_{i=1}^{n}\mu_{i}}{r}\,.
    \tag{3.21}
  \end{equation}
  若 $D^{r}_{g;\mu}$ 不是整数，相应的 Witten 类为零。$r=2$ 的情形给出平凡上同调场论：$W^{2}_{g;1,\dots,1}=\mathbf{1}_{g,n}$。在亏格零情形，这一构造首先由 Witten [Wit93] 借助 $r$-旋结构给出。Witten 类在高亏格情形的构造，首先由 Polishchuk 与 Vaintrob [PV00] 得到。相应的配分函数是一个 $r$-KdV tau 函数 [FSZ10]，并且它满足 W-约束，这等价于在 $r$-Airy 谱曲线 $\frac{1}{r}y^{r}-x=0$ 上的拓扑递归。
% [Tutor Note: 术语说明：原文 $r$-spin class 这里译作“$r$-旋类”，spin structure 译作“旋结构”（代数几何惯例）；若全书统一改为“自旋”，需全文一并替换。]
% [Tutor Note: (3.21) 的核对：令 $r=2$，则分子为 $(0)(g-1)-n+\sum_{i}\mu_{i}=\sum_{i}\mu_{i}-n$；而 $r=2$ 时 $V=\Q\cdot e_{1}$ 是维数 $1$ 的，所有 $\mu_{i}=1$，故 $D^{2}_{g;\mu}=0$，与原文“$r=2$ 给出平凡上同调场论 $W^{2}_{g;1,\dots,1}=\mathbf{1}_{g,n}$”完全一致。另注意 (3.20) 的指数 $2D^{r}_{g;\mu}$ 是实次数，故 $D^{r}_{g;\mu}$ 为半整数时该上同调群为零。]
  \item 在 [PPZ15] 中证明了，$\Mc_{g,n}$ 的所谓重言环（即 $\Mc_{g,n}$ 上同调中对重言映射的前推与拉回封闭的最小子代数）里所有已知的关系，都可由 Witten $r$-旋类导出。
% [Tutor Note: 配对 $\eta(e_{\mu},e_{\nu})=\delta_{\mu+\nu,r}$ 是非退化的：$e_{\mu}$ 只与 $e_{r-\mu}$ 配对，而 $1\le r-\mu\le r-1$，故这是 $V$ 上的一个完美配对（无需再用 $\eta^{\mu\nu}$ 的矩阵求逆即可看出）。核对：$r=2$ 时 $\eta(e_{1},e_{1})=\delta_{2,2}=1$，正与 Example 3.2 的一维约定吻合。]
  \item 在 [CGG25] 中，作者引入了 $\Theta$ 类的 $r$-旋版本，记作 $\Theta^{r}_{g,n}$，它满足与 Witten 类相类似的性质。该文证明了它满足 W-约束，这等价于在 $r$-Bessel 谱曲线 $\frac{1}{r}y^{r}x-1=0$ 上的拓扑递归。依赖于第二个整数 $s\in\{1,\dots,r-1\}$ 的进一步推广 $\Theta^{r,s}_{g,n}$ 在 [BCGS] 中得到研究，那里证明了 W-约束以及一个推广版本的拓扑递归。
  \item 在 [Chio8] 中，Chiodo 定义了 Hodge 类的一个推广，称为\emph{$\Omega$ 类}，它依赖两个整数 $r\ge 1$ 与 $s\in\Z$。取 $V=\bigoplus_{\mu=1}^{r}\Q\cdot e_{\mu}$，配对由 $\eta(e_{\mu},e_{\nu})=\frac{1}{r}\delta_{\mu+\nu\equiv 0\ (\mathrm{mod}\ r)}$ 给出。$\Omega$ 类是一个具有混合上同调次数的 CohFT：
  \begin{equation}
    \Omega^{r,s}_{g;\mu_{1},\dots,\mu_{n}}\in H^{2\bullet}(\Mc_{g,n},\Q)\,.
    \tag{3.22}
  \end{equation}
  它定义为黎曼曲面模空间上一个（虚拟）向量丛的全 Chern 类。$r=s=1$ 的情形还原 Hodge 类。$r\ge 2$ 且 $s=\pm1$ 的情形分别与 Witten $r$-旋类和 Theta $r$-旋类有关。
% [Tutor Note: Chiodo 配对中的下标是一个同余条件，而非 Kronecker 的 $\mu+\nu=r$：$\eta(e_{\mu},e_{\nu})=\frac{1}{r}$ 当且仅当 $\mu+\nu\equiv 0\ (\mathrm{mod}\ r)$，否则为 $0$。核对：$r=1$ 时 $V=\Q\cdot e_{1}$ 且 $\eta(e_{1},e_{1})=1$，退回 Hodge 类的一维情形；$r=2$ 时 $\eta(e_{1},e_{1})=\eta(e_{2},e_{2})=\frac{1}{2}$、$\eta(e_{1},e_{2})=0$，仍非退化。]
  \item 设 $G$ 是一个复、单、单连通的 Lie 群，Lie 代数为 $\mathfrak{g}$。固定整数 $\ell>0$，令 $V$ 为由 $\mathfrak{g}$ 在能级 $\ell$ 处的不可约表示 $e_{\mu}$ 张成的 $\Q$-向量空间。令 $\eta(e_{\mu},e_{\nu})=\delta_{\mu,\nu^{\star}}$，其中 $\nu^{\star}$ 表示对偶表示，并令 $e_{1}$ 为与平凡表示对应的向量。Verlinde 丛是向量丛
  \begin{equation}
    \mathcal{V}^{g,\ell}_{g;\mu_{1},\dots,\mu_{n}}\longrightarrow\Mc_{g,n}
    \tag{3.23}
  \end{equation}
  它在光滑黎曼曲面上的纤维是非阿贝尔 theta 函数空间。Verlinde 丛的 Chern 特征
  \begin{equation}
    \mathcal{V}^{g,\ell}_{g;\mu_{1},\dots,\mu_{n}}:=\mathrm{ch}\big(\mathcal{V}^{g,\ell}_{g;\mu_{1},\dots,\mu_{n}}\big)\in H^{2\bullet}(\Mc_{g,n},\Q)
    \tag{3.24}
  \end{equation}
  构成一个带单位的 CohFT [MOPPZ17]。其粘合公理是融合规则（fusion rules）的推论，而单位公理则是真空的传播（propagation of vacua）。
\end{itemize}
\end{example}

% [Tutor Note: (3.23)-(3.24) 的记号有两处值得注意（只提示，不改原文）：其一，(3.24) 用同一个符号 $\mathcal{V}^{g,\ell}_{g;\mu_{1},\dots,\mu_{n}}$ 既表示丛（$\mathrm{ch}$ 的辐角）又表示它的 Chern 特征（等号左端），严格写法应对两者加以区分（例如丛写作 $\mathbb{V}$）；其二，上标 $g,\ell$ 中的 $g$ 与下标中的 $g$ 重复，文献中通常只标 level $\ell$。]
% [Tutor Note: Example 3.2 的最后一句（“Verlinde 丛的 Chern 特征 (3.24) 构成一个带单位的 CohFT [MOPPZ17]……”）在原文中跨 p.30-p.31，续文在 p.31 顶部。按“跨页句子由开始者写完”的规定，本块已写完该句并在 (3.24) 之后闭合 Example 3.2。]
% [原文 p.31 续句见上：Block G 应从 p.31 的第二条（公式 (3.25) 所在的“固定目标 Kähler 时空”一条）接续。]
